%! Representing a Categorical Variable — Unit 1, Lesson 1.2 — AP Practice (Answer Key)
% GRADUAL RELEASE: AP-format practice, assigned but NOT scored — the cover's score column
% reads NA for this component. The graded practice is the `homework` component that follows it
% at the back of the packet.
% This page also CARRIES THE TRANSFER: the recycling audit and the activity survey appear
% nowhere else in the packet, so these are the first time students meet either data set.
% Part (a) of the free response is the spiral item — it reaches back to Lesson 1.1 and makes
% students classify the variable before they are allowed to summarize it.
% GENERATED FROM homework/main.tex. Each \writelines{n} in the blank is answered by an
% \ansline{} terse enough not to wrap past n lines — that is what keeps the page counts equal.
% NO teachernote here — scoring guidance lives in the lesson plan.
\documentclass[10pt]{article}
\usepackage{apstats-article}
\usepackage{apstats-key}   % pulls in -boxes; do NOT also load -boxes

\begin{document}

\pageheader{Unit 1, Lesson 1.2}{AP Practice --- Answer Key}
% NAMESTRIP: no \namedateperiod here — mirrors the blank.

\vspace{0.15in}

% Authored BYTE-IDENTICALLY in ap_practice/ and ap_practice_key/ — this box is what keeps the
% blank and the key the same number of pages.
\boxguard[8]
\begin{remindbox}
\small
This page is \textbf{not required and not scored} --- your graded homework is the last section
of this packet. These are AP-format reps: work them for the practice, and bring what you get
stuck on to office hours or study hall.
\end{remindbox}

\vspace{0.12in}

\boxguard[26]
\begin{scenariobox}[The Recycling Audit]{navy}
\small
A city inspected recycling bins in two neighborhoods and recorded what each bin mainly
contained. Riverside had \textbf{250} bins inspected; Hillcrest had \textbf{500}.

\vspace{8pt}
\renewcommand{\arraystretch}{1.25}
\begin{tabularx}{\linewidth}{@{} l Y Y Y Y Y @{}}
\rowcolor{sky}
\textbf{Neighborhood} & \textbf{Paper} & \textbf{Plastic} & \textbf{Glass} &
\textbf{Contaminated} & \textbf{Total} \\
\hline
Riverside & 100 &  75 &  50 & 25 & 250 \\
Hillcrest & 190 & 160 & 110 & 40 & 500 \\
\hline
\end{tabularx}
\end{scenariobox}

\vspace{0.14in}

\boxguard[16]
\begin{headlinebox}{goldbg}
\small
\textbf{Section I --- Multiple Choice.} Circle the single best answer. Every answer choice is
about the context, not just the arithmetic --- read them all before choosing.
\end{headlinebox}

\vspace{0.10in}

\begin{enumerate}[leftmargin=1.9em, itemsep=0.13in]

  \item What percent of the Riverside bins inspected were contaminated?
        \begin{enumerate}[label=\textbf{(\Alph*)}, leftmargin=2.2em, itemsep=2pt]
          \item 2.5\%
          \item 8\%
          \item \ans{10\%}
          \item 25\%
          \item 40\%
        \end{enumerate}

  \item Which statement about contamination in the two neighborhoods is \textbf{correct}?
        \begin{enumerate}[label=\textbf{(\Alph*)}, leftmargin=2.2em, itemsep=2pt]
          \item Hillcrest is worse, because it had more contaminated bins.
          \item Riverside is worse, because it had more contaminated bins.
          \item \ans{Hillcrest had more contaminated bins, but Riverside had the higher
                contamination rate.}
          \item Riverside had more contaminated bins, but Hillcrest had the higher
                contamination rate.
          \item The neighborhoods cannot be compared, because different numbers of bins
                were inspected.
        \end{enumerate}

  \item A relative frequency table is built for Riverside's four categories. Which must
        be true?
        \begin{enumerate}[label=\textbf{(\Alph*)}, leftmargin=2.2em, itemsep=2pt]
          \item The four entries add to 250.
          \item \ans{The four entries add to 1, or equivalently 100\%.}
          \item The four entries add to 750, the total bins inspected in the city.
          \item Each entry must be larger than the matching entry for Hillcrest.
          \item The entries cannot be compared to Hillcrest's under any circumstances.
        \end{enumerate}

  \item The city wants one graph showing whether the two neighborhoods sort their recycling
        differently. Which display best supports that comparison?
        \begin{enumerate}[label=\textbf{(\Alph*)}, leftmargin=2.2em, itemsep=2pt]
          \item Side-by-side bars of the raw counts, because counts are the actual data.
          \item \ans{Side-by-side bars of the relative frequencies, because the neighborhoods had
                different numbers of bins.}
          \item A single bar for each neighborhood's total.
          \item Bars of the raw counts with Hillcrest's halved to match Riverside's total.
          \item No graph can compare two neighborhoods.
        \end{enumerate}

\end{enumerate}

\vspace{0.14in}

\boxguard[16]
\begin{headlinebox}{goldbg}
\small
\textbf{Section II --- Free Response.} Show your reasoning. Every answer must be written
\emph{in the context of the problem} --- that is what earns credit on the AP exam.
\end{headlinebox}

\vspace{0.10in}

\begin{enumerate}[leftmargin=1.9em, itemsep=0.16in, resume]

  \item A school surveyed its 9th graders and its 12th graders about how they would prefer to
        spend an activity period. Among \textbf{300} ninth graders: Sports 150, Clubs 90,
        Study hall 60. Among \textbf{120} twelfth graders: Sports 42, Clubs 30, Study hall 48.

        \vspace{6pt}
        \textbf{(a)} Identify the individuals and the variable, and classify the variable as
        categorical or quantitative.\\[4pt]
        \ansline{Individuals: the 300 ninth graders and 120 twelfth graders surveyed. The}\\[1pt]
        \ansline{variable is preferred use of an activity period, and it is categorical.}

        \vspace{6pt}
        \textbf{(b)} Build the relative frequency table for the twelfth graders. Show each
        calculation.\\[4pt]
        \ansline{Sports $42/120 = 0.35 = 35\%$; Clubs $30/120 = 0.25 = 25\%$;}\\[1pt]
        \ansline{Study hall $48/120 = 0.40 = 40\%$. The three add to $1.00$, as they must,}\\[1pt]
        \ansline{since every senior surveyed picked exactly one option.}

        \vspace{6pt}
        \textbf{(c)} A student newspaper writes: \textbf{``Ninth graders like sports far more
        than seniors do --- 150 of them chose it, versus only 42 seniors.''} Explain what is
        wrong with this comparison, and give the corrected version.\\[4pt]
        \ansline{The grades are different sizes --- 300 ninth graders against 120 seniors ---}\\[1pt]
        \ansline{so raw counts are not comparable. Corrected: 50\% of ninth graders chose}\\[1pt]
        \ansline{sports versus 35\% of seniors, a real but far smaller gap than 150 to 42.}

        \vspace{6pt}
        \textbf{(d)} Identify one preference where the two grades differ most, and describe
        that difference in context using relative frequencies.\\[4pt]
        \ansline{Study hall differs most: 20\% of ninth graders chose it against 40\% of}\\[1pt]
        \ansline{seniors --- seniors picked study hall at twice the rate of ninth graders.}

\end{enumerate}

\vspace{0.12in}

\boxguard[12]
\begin{extensionbox}
\small
Find a news story that compares two groups using raw counts --- two cities, two countries, two
schools. Write down the counts, then look up the sizes of the two groups and redo the
comparison as rates. Say whether the story's point survives.\\[4pt]
\ansline{Answers vary. Look for a story where the larger group leads on counts;}\\[1pt]
\ansline{converting to rates often reverses or erases the reported gap.}
\end{extensionbox}

\vspace{0.12in}


\end{document}
